\section*{Chapter \ref{ch:qm:jump-processes} --- Jump Processes}

\begin{solution}{exo:qm:jump-processes:1}
$e^{-5/12} = 65.9\%$; by memorylessness the expected wait is $1/5$ of a year, whatever has happened.
\end{solution}

\begin{solution}{exo:qm:jump-processes:2}
Mean $3 \times (-2\%) = -6\%$; variance $3(0.02^2 + 0.04^2) = 0.006$, a standard deviation of 7.75\%.
\end{solution}

\begin{solution}{exo:qm:jump-processes:3}
For $s < t$, $N_t - N_s$ is independent of $\mathcal F_s$ with mean $\lambda(t - s)$, so $\E_s[N_t - \lambda t] = N_s - \lambda s$;
and $\E[(N_t - \lambda t)^2] = \Var(N_t) = \lambda t$.
\end{solution}

\begin{solution}{exo:qm:jump-processes:4}
$d(e^{X_t}) = e^{X_{t-}}(\sigma\,dW_t + \tfrac12\sigma^2dt) + e^{X_{t-}}(e^{\Delta X_t} - 1)$. The jump term has compensator
$e^{X_{t-}}\lambda(\E[e^Y] - 1)\,dt$, so the drift $b = -\tfrac12\sigma^2 - \lambda(\E[e^Y] - 1)$ added to $X$ makes $e^X$ a
\omterm{def:qm:probability-at-speed:martingale}{martingale}.
\end{solution}

\begin{solution}{exo:qm:jump-processes:5}
$\psi(u) = 2(50/(50 - \iu u) - 1)$; mean $\lambda\E[Y] = 2/50 = 0.04$, variance $\lambda\E[Y^2] = 2 \times 2/50^2 = 0.0016$.
\end{solution}

\begin{solution}{exo:qm:jump-processes:6}
$76.4/5 = 15.3$ at a week and $76.4/252 = 0.30$ at a year.
\end{solution}

\begin{solution}{exo:qm:jump-processes:7}
$(-0.1,\ 0.043,\ -0.00378,\ 0.001888)$ both ways: the finite differences of the \omterm{def:qm:jump-processes:cumulant}{cumulant} generating
function agree with the formulas to about $10^{-7}$.
\end{solution}

\begin{solution}{exo:qm:jump-processes:8}
In a Lévy model the excess kurtosis falls like $1/t$: 1.2 at 21 days corresponds to about 25 at one day.
Monthly returns average away exactly the jumps that matter for one-day risk; measure the tail at the
horizon of the risk.
\end{solution}

\begin{solution}{pb:qm:jump-processes:1}
\textbf{1.} $20.98 \times 0.98\% = 20.6\%$.
\textbf{2.} $\log_{10}\Phi(-20.98) = -97.3$.
\textbf{3.} $20\%/0.98\% = 20.4$ standard deviations: $\log_{10}$ probability $-92.2$.
\textbf{4.} $0.98\%\sqrt{252} = 15.6\%$.
\textbf{5.} A Gaussian volatility high enough to make the day plausible would make ordinary days
several times more volatile than they are; the tail and the body cannot both fit.
\textbf{6.} $\mu_J = -11.1\%$.
\textbf{7.} 13.0\% a year (12.96\%).
\textbf{8.} 31\%: $0.5/252 \times (0.111^2 + 0.05^2)$ out of $0.0098^2$.
\textbf{9.} $1 - e^{-0.5/252} = 0.20\%$.
\textbf{10.} With probability 0.116\% a day: once every 3.4 years.
\textbf{11.} Skewness $-4.64$, excess kurtosis 76.4.
\textbf{12.} $76.4/21 = 3.64$.
\textbf{13.} \omterm{def:qm:jump-processes:cumulant}{Cumulants} of a \omterm{def:qm:jump-processes:levy}{Lévy process} are linear in $t$, so $\kappa_4/\kappa_2^2 \propto t/t^2$.
\textbf{14.} $-\lambda(\E[e^Y] - 1) = -0.5(e^{-0.111 + 0.00125} - 1) = +5.2\%$ a year, on top of $-\tfrac12\sigma_c^2$.
\textbf{15.} 15.3 in theory, 15.6 in 400\,000 simulated five-day returns.
\textbf{16.} A choice: one such day in the record cannot estimate a frequency; the model encodes a
judgement about how often it should be allowed to happen.
\textbf{17.} Volatility clustering: in the data the kurtosis decays more slowly than $1/t$ (chapter~18).
\textbf{18.} Short-dated, far out-of-the-money puts, whose value is almost all jump probability (One
Quant Book~5, chapter~13).
\textbf{19.} Named result: \emph{twenty standard deviations}: the 1987 day was a 21-standard-deviation
fall, with Gaussian probability $10^{-97}$; jumps at rate 0.5 a year with mean $-11.1\%$ and standard
deviation 5\%, carrying 31\% of the variance, make a fall of 20\% a once-in-fifty-years event, at the
cost of a one-day excess kurtosis of 76 (3.6 at a month).
\textbf{20.} It changes the shape of the one-day law, putting mass in the tail without raising the
variance of ordinary days.
\end{solution}

\begin{solution}{iq:qm:jump-processes:1}
That the model, not the market, is wrong: under it the move should never have happened. Check the data
first (a bad print, a corporate action), then the model's assumptions (continuous paths, constant
volatility, Gaussian shocks) and the risk numbers that depend on them.
\iqlookfor{model failure, not bad luck.}
\end{solution}

\begin{solution}{iq:qm:jump-processes:2}
$1 - e^{-2}(1 + 2) = 59.4\%$; half a second, whatever the time since the last trade.
\iqlookfor{Poisson probabilities and memorylessness.}
\end{solution}

\begin{solution}{iq:qm:jump-processes:3}
A delta hedge offsets small moves to first order; a jump is a large move, and the option's value changes
by the full difference $V(S + \Delta S) - V(S)$, not $\Delta\cdot\Delta S$. With random jump sizes there is one more
source of risk than the stock can hedge: the market is incomplete.
\iqlookfor{nonlinearity plus incompleteness.}
\end{solution}

\begin{solution}{iq:qm:jump-processes:4}
$\psi(u) = \iu\gamma u - \tfrac12\sigma^2u^2 + \int(e^{\iu ux} - 1 - \iu ux\mathbf 1_{|x|<1})\,\nu(dx)$: $\gamma$ a drift, $\sigma^2$ the
Brownian variance, $\nu$ the expected number of jumps per unit time by size.
\iqlookfor{the three parts and the role of $\nu$.}
\end{solution}

\begin{solution}{iq:qm:jump-processes:5}
\omterm{def:qm:jump-processes:cumulant}{Cumulants} scale with $t$, so the excess kurtosis falls like $1/t$ and a monthly return is nearly
Gaussian. In the data it falls more slowly, because volatility clusters and increments are not
independent.
\iqlookfor{$\kappa_n(t) = t\kappa_n(1)$ and what breaks it.}
\end{solution}

\begin{solution}{iq:qm:jump-processes:6}
The discounted price must be a \omterm{def:qm:probability-at-speed:martingale}{martingale}: for $S = S_0e^X$ with $X = bt + \sigma W + J$, $b = r - \tfrac12\sigma^2 -
\lambda(\E[e^Y] - 1)$, the jump compensator removing the jumps' average contribution. Otherwise the model
admits arbitrage against the \omterm{def:qm:girsanov-and-changes-of-numeraire:numeraire}{money-market account}.
\iqlookfor{the compensator in the drift.}
\end{solution}
